Nesse apêndice estão contidas os dados referentes as respostas de maior
relevância para a implementação do sistema. Considerando principalmente
que ela possui o total de vinte perguntas.

\begin{table}[ht]
  \centering
  \begin{tabular}{
    |p{\dimexpr0.8\textwidth-4.8\tabcolsep-3.2\arrayrulewidth\relax}|
    >{\raggedleft\arraybackslash}p{\dimexpr0.1\textwidth-0.6\tabcolsep-0.4\arrayrulewidth\relax}|
    >{\raggedleft\arraybackslash}p{\dimexpr0.1\textwidth-0.6\tabcolsep-0.4\arrayrulewidth\relax}|
  }
    \hline
    \multicolumn{3}{
      |>{\centering\arraybackslash}
      p{\dimexpr\textwidth-2\tabcolsep-2\arrayrulewidth\relax}|
    }{
      \textbf{
        1. Você já teve algum interesse ou experiência com o monitoramento
        do seu humor ou estado emocional?
      }
    } \\
    \hline
    \multicolumn{1}{|c|}{} &
    \multicolumn{1}{c|}{\textbf{\%}} &
    \multicolumn{1}{c|}{\textbf{TOTAL}} \\
    \hline
    Sim, faço regularmente & 18,8 & 9 \\
    \hline
    Sim, já tentei alguma vez & 20,8 & 10 \\
    \hline
    Tenho interesse, mas nunca tentei & 52,1 & 25 \\
    \hline
    Não tenho interesse e nunca tentei & 8,3 & 4 \\
    \hline
  \end{tabular}
  \source{Elaborado pelo autor (2026)}
\end{table}


\begin{table}[ht]
  \centering
  \begin{tabular}{
    |p{\dimexpr0.8\textwidth-4.8\tabcolsep-3.2\arrayrulewidth\relax}|
    >{\raggedleft\arraybackslash}p{\dimexpr0.1\textwidth-0.6\tabcolsep-0.4\arrayrulewidth\relax}|
    >{\raggedleft\arraybackslash}p{\dimexpr0.1\textwidth-0.6\tabcolsep-0.4\arrayrulewidth\relax}|
  }
    \hline
    \multicolumn{3}{
      |>{\centering\arraybackslash}
      p{\dimexpr\textwidth-2\tabcolsep-2\arrayrulewidth\relax}|
    }{
      \textbf{
        2. Respostas sobre a faixa etária dos participantes
      }
    } \\
    \hline
    \multicolumn{1}{|c|}{} &
    \multicolumn{1}{c|}{\textbf{\%}} &
    \multicolumn{1}{c|}{\textbf{TOTAL}} \\
    \hline
    Menos de 18 anos & 0 & 0 \\
    \hline
    Entre 18 e 30 anos & 43,8 & 21 \\
    \hline
    31 a 45 anos & 29,2 & 14 \\
    \hline
    46 a 60 anos & 22,9 & 11 \\
    \hline
    Mais de 60 anos & 4,2 & 2 \\
    \hline
  \end{tabular}
  \source{Elaborado pelo autor (2026)}
\end{table}


\begin{table}[ht]
  \centering
  \begin{tabular}{
    |p{\dimexpr0.8\textwidth-4.8\tabcolsep-3.2\arrayrulewidth\relax}|
    >{\raggedleft\arraybackslash}p{\dimexpr0.1\textwidth-0.6\tabcolsep-0.4\arrayrulewidth\relax}|
    >{\raggedleft\arraybackslash}p{\dimexpr0.1\textwidth-0.6\tabcolsep-0.4\arrayrulewidth\relax}|
  }
    \hline
    \multicolumn{3}{
      |>{\centering\arraybackslash}
      p{\dimexpr\textwidth-2\tabcolsep-2\arrayrulewidth\relax}|
    }{
      \textbf{
        3. Com qual gênero você se identifica?
      }
    } \\
    \hline
    \multicolumn{1}{|c|}{} &
    \multicolumn{1}{c|}{\textbf{\%}} &
    \multicolumn{1}{c|}{\textbf{TOTAL}} \\
    \hline
    Feminino & 66,7 & 32 \\
    \hline
    Masculino & 33,3 & 16 \\
    \hline
    Não-binário & 0 & 0 \\
    \hline
    Prefiro não dizer & 0 & 0 \\
    \hline
  \end{tabular}
  \source{Elaborado pelo autor (2026)}
\end{table}

\begin{table}[ht]
  \centering
  \begin{tabular}{
    |p{\dimexpr0.8\textwidth-4.8\tabcolsep-3.2\arrayrulewidth\relax}|
    >{\raggedleft\arraybackslash}p{\dimexpr0.1\textwidth-0.6\tabcolsep-0.4\arrayrulewidth\relax}|
    >{\raggedleft\arraybackslash}p{\dimexpr0.1\textwidth-0.6\tabcolsep-0.4\arrayrulewidth\relax}|
  }
    \hline
    \multicolumn{3}{
      |>{\centering\arraybackslash}
      p{\dimexpr\textwidth-2\tabcolsep-2\arrayrulewidth\relax}|
    }{
      \textbf{
        3. Você já recebeu algum diagnóstico relacionado à saúde mental?
      }
    } \\
    \hline
    \multicolumn{1}{|c|}{} &
    \multicolumn{1}{c|}{\textbf{\%}} &
    \multicolumn{1}{c|}{\textbf{TOTAL}} \\
    \hline
    Sim & 62,5 & 30 \\
    \hline
    Não & 37,5 & 18 \\
    \hline
    Prefiro não dizer & 0 & 0 \\
    \hline
  \end{tabular}
  \source{Elaborado pelo autor (2026)}
\end{table}

\begin{table}[ht]
  \centering
  \begin{tabular}{
    |p{\dimexpr0.8\textwidth-4.8\tabcolsep-3.2\arrayrulewidth\relax}|
    >{\raggedleft\arraybackslash}p{\dimexpr0.1\textwidth-0.6\tabcolsep-0.4\arrayrulewidth\relax}|
    >{\raggedleft\arraybackslash}p{\dimexpr0.1\textwidth-0.6\tabcolsep-0.4\arrayrulewidth\relax}|
  }
    \hline
    \multicolumn{3}{
      |>{\centering\arraybackslash}
      p{\dimexpr\textwidth-2\tabcolsep-2\arrayrulewidth\relax}|
    }{
      \textbf{
        Você realiza algum acompanhamento psicológico ou psiquiátrico
        atualmente?
      }
    } \\
    \hline
    \multicolumn{1}{|c|}{} &
    \multicolumn{1}{c|}{\textbf{\%}} &
    \multicolumn{1}{c|}{\textbf{TOTAL}} \\
    \hline
    Sim, com psicólogo(a) & 12,5 & 6 \\
    \hline
    Sim, com psiquiatra & 4,2 & 2 \\
    \hline
    Sim, com ambos & 12,5 & 6 \\
    \hline
    Não & 70,8 & 34 \\
    \hline
  \end{tabular}
  \source{Elaborado pelo autor (2026)}
\end{table}

\begin{table}[ht]
  \centering
  \begin{tabular}{
    |p{\dimexpr0.8\textwidth-4.8\tabcolsep-3.2\arrayrulewidth\relax}|
    >{\raggedleft\arraybackslash}p{\dimexpr0.1\textwidth-0.6\tabcolsep-0.4\arrayrulewidth\relax}|
    >{\raggedleft\arraybackslash}p{\dimexpr0.1\textwidth-0.6\tabcolsep-0.4\arrayrulewidth\relax}|
  }
    \hline
    \multicolumn{3}{
      |>{\centering\arraybackslash}
      p{\dimexpr\textwidth-2\tabcolsep-2\arrayrulewidth\relax}|
    }{
      \textbf{
        Você já monitorou seu humor em algum momento da vida?
      }
    } \\
    \hline
    \multicolumn{1}{|c|}{} &
    \multicolumn{1}{c|}{\textbf{\%}} &
    \multicolumn{1}{c|}{\textbf{TOTAL}} \\
    \hline
    Sim, regularmente & 21,3 & 10 \\
    \hline
    Sim, de vez em quando & 12,8 & 6 \\
    \hline
    Já tentei, mas parei & 21,3 & 10 \\
    \hline
    Nunca & 44,7 & 21 \\
    \hline
  \end{tabular}
\end{table}


\begin{table}[ht]
  \centering
  \begin{tabular}{
    |p{\dimexpr0.8\textwidth-4.8\tabcolsep-3.2\arrayrulewidth\relax}|
    >{\raggedleft\arraybackslash}p{\dimexpr0.1\textwidth-0.6\tabcolsep-0.4\arrayrulewidth\relax}|
    >{\raggedleft\arraybackslash}p{\dimexpr0.1\textwidth-0.6\tabcolsep-0.4\arrayrulewidth\relax}|
  }
    \hline
    \multicolumn{3}{
      |>{\centering\arraybackslash}
      p{\dimexpr\textwidth-2\tabcolsep-2\arrayrulewidth\relax}|
    }{
      \textbf{
        Como você costuma ou costumava fazer esse monitoramento?
      }\footnotemark
    } \\
    \hline
    \multicolumn{1}{|c|}{} &
    \multicolumn{1}{c|}{\textbf{\%}} &
    \multicolumn{1}{c|}{\textbf{TOTAL}} \\
    \hline
    Apenas mentalmente & 48,3 & 14 \\
    \hline
    Diário escrito & 37,9 & 11 \\
    \hline
    Aplicativos & 34,5 & 10 \\
    \hline
    Afetivograma & 10,3 & 3 \\
    \hline
    Planilhas & 6,9 & 2 \\
    \hline
    Auto observação & 3,4 & 1 \\
    \hline
  \end{tabular}
\end{table}
\footnotetext{Questão de múltipla escolha (respostas não somam 100\%). Base: 29 respondentes que relataram ter monitorado o humor de alguma forma.}

\begin{table}[ht]
  \centering
  \begin{tabular}{
    |p{\dimexpr0.8\textwidth-4.8\tabcolsep-3.2\arrayrulewidth\relax}|
    >{\raggedleft\arraybackslash}p{\dimexpr0.1\textwidth-0.6\tabcolsep-0.4\arrayrulewidth\relax}|
    >{\raggedleft\arraybackslash}p{\dimexpr0.1\textwidth-0.6\tabcolsep-0.4\arrayrulewidth\relax}|
  }
    \hline
    \multicolumn{3}{
      |>{\centering\arraybackslash}
      p{\dimexpr\textwidth-2\tabcolsep-2\arrayrulewidth\relax}|
    }{
      \textbf{
        Se você parou ou teve dificuldades, quais foram os motivos?
      }
    } \\
    \hline
    \multicolumn{1}{|c|}{} &
    \multicolumn{1}{c|}{\textbf{\%}} &
    \multicolumn{1}{c|}{\textbf{TOTAL}} \\
    \hline
    Esquecimento & 56,0 & 14 \\
    \hline
    Falta de tempo & 40,0 & 10 \\
    \hline
    Preocupações com privacidade & 24,0 & 6 \\
    \hline
    Não percebi benefícios & 20,0 & 5 \\
    \hline
    Ferramentas difíceis de usar & 8,0 & 2 \\
    \hline
    Dificuldade de analisar sozinho & 4,0 & 1 \\
    \hline
  \end{tabular}
\end{table}

\begin{table}[ht]
  \centering
  \begin{tabular}{
    |p{\dimexpr0.8\textwidth-4.8\tabcolsep-3.2\arrayrulewidth\relax}|
    >{\raggedleft\arraybackslash}p{\dimexpr0.1\textwidth-0.6\tabcolsep-0.4\arrayrulewidth\relax}|
    >{\raggedleft\arraybackslash}p{\dimexpr0.1\textwidth-0.6\tabcolsep-0.4\arrayrulewidth\relax}|
  }
    \hline
    \multicolumn{3}{
      |>{\centering\arraybackslash}
      p{\dimexpr\textwidth-2\tabcolsep-2\arrayrulewidth\relax}|
    }{
      \textbf{
        Quão importante é para você ter controle total sobre seus dados pessoais no aplicativo?
      }
    } \\
    \hline
    \multicolumn{1}{|c|}{} &
    \multicolumn{1}{c|}{\textbf{\%}} &
    \multicolumn{1}{c|}{\textbf{TOTAL}} \\
    \hline
    Muito importante & 68,1 & 32 \\
    \hline
    Importante & 12,8 & 6 \\
    \hline
    Neutro & 12,8 & 6 \\
    \hline
    Pouco importante & 2,1 & 1 \\
    \hline
    Nada importante & 4,3 & 2 \\
    \hline
  \end{tabular}
\end{table}

\begin{table}[ht]
  \centering
  \begin{tabular}{
    |p{\dimexpr0.8\textwidth-4.8\tabcolsep-3.2\arrayrulewidth\relax}|
    >{\raggedleft\arraybackslash}p{\dimexpr0.1\textwidth-0.6\tabcolsep-0.4\arrayrulewidth\relax}|
    >{\raggedleft\arraybackslash}p{\dimexpr0.1\textwidth-0.6\tabcolsep-0.4\arrayrulewidth\relax}|
  }
    \hline
    \multicolumn{3}{
      |>{\centering\arraybackslash}
      p{\dimexpr\textwidth-2\tabcolsep-2\arrayrulewidth\relax}|
    }{
      \textbf{
        Que tipo de tom você preferiria em um aplicativo de humor?
      }
    } \\
    \hline
    \multicolumn{1}{|c|}{} &
    \multicolumn{1}{c|}{\textbf{\%}} &
    \multicolumn{1}{c|}{\textbf{TOTAL}} \\
    \hline
    Acolhedor e empático & 37,5 & 18 \\
    \hline
    Neutro e objetivo & 25,0 & 12 \\
    \hline
    Clínico e profissional & 20,8 & 10 \\
    \hline
    Lúdico e leve & 12,5 & 6 \\
    \hline
    Moderno e minimalista & 2,1 & 1 \\
    \hline
    Não gosto de apps que parecem um formulário psiquiátrico & 2,1 & 1 \\
    \hline
  \end{tabular}
\end{table}

\FloatBarrier

\newpage
Para além dos itens citados acima, no formulário existia a opção de
deixar uma contribuição, essa lista apresenta as sugestões de
funcionalidades adicionais obtidas a partir das respostas dos
participantes na questão de entrada livre:

\begin{itemize}
  \item Questões hormonais e voltadas para mulheres que acompanham as
  mudanças de humor no ciclo menstrual.
  \item Seria legal ver um gráfico semanal ou mensal do meu humor.
  \item Algo totalmente offline ou que eu possa apagar quando quiser
  sem me preocupar.
  \item Sugestões de atividades ou conteúdos com base no humor
  registrado (música, respiração guiada, frases motivacionais, pequenas metas diárias).
  \item Design que considere a acessibilidade para pessoas neuro
  divergentes (cores e tipografia suaves).
  \item Opção de visualizar gráficos ao longo do tempo.
  \item Bloco de notas para pensamentos, reflexões e desabafos.
  \item Função para considerar o ciclo menstrual no humor.
  \item Possibilidade de escrita, registro de sintomas físicos,
  sugestões de regulação.
\end{itemize}

Devido à natureza livre das respostas, esses dados foram normalizados
e removidos duplicatas ou elementos julgados nulos
(respostas como ``nada a acrescentar'', etc).
